% GENERATED FILE. Edit canonical lessons, then run npm run generate:books.
% canonical-source-hash: fnv1a64:90dd2787819bab86
% canonical-lessons: GE-C334-zustimmen, GE-C334-zustimmung, GE-C334-einverstanden, GE-C334-zugeben, GE-C334-einraeumen

\chapter{Zustimmen und zugeben}
\label{ch:ge-zustimmen-und-zugeben}
\hlchaptermodality{334}

\begin{chapteropening}
\textbf{By the end of this chapter:} \emph{I can agree with someone, say I am happy with a plan, admit a mistake and concede a point.}

Complete the last lesson of chapter 334: I can agree with someone, say I am happy with a plan, admit a mistake and concede a point.
\end{chapteropening}

\section[zustimmen]{zustimmen — to agree (with)}

\label{lesson:GE-C334-zustimmen}



\begin{quote}
Before the new one: say the German for to object, then the German for to deny.
\end{quote}

\subsection*{You'll want to know: zustimmen}
\textbf{zustimmen} — \textquotedblleft{}to agree (with)\textquotedblright{}.

\begin{quote}
Ich stimme dir zu: Wir sollten früher losfahren.
\end{quote}

\begin{quote}
\emph{I agree with you: we should set off earlier.}
\end{quote}

Separable, with the dative: \emph{ich stimme dir zu}, \emph{sie hat zugestimmt}. \emph{Einem Plan zustimmen} — to approve a plan.

\begin{cousinweb}
\emph{zustimmen} is built from \emph{zu} \textquotedblleft{}to\textquotedblright{} and \emph{stimmen}, a verb made from \emph{Stimme} \textquotedblleft{}voice, vote\textquotedblright{} (EtymDB): you add your voice to an idea.
\end{cousinweb}

\begin{grammarlens}[title={What you've built}]
One more word for telling, explaining and describing at B1.
\end{grammarlens}

\subsection*{Your turn}
\begin{itemize}
\raggedright
  \item \emph{Say it:} \emph{zustimmen}
  \item \emph{Say it:} \emph{zustimmen}, once more
  \item \emph{Say it:} a sentence of your own with \emph{zustimmen}
  \item \emph{Recall:} say the German for to object, then the German for to deny, then say \emph{zustimmen} again
\end{itemize}

\begin{tcolorbox}[breakable,colback=teal!4,colframe=teal!35!black,title={Before you move on}]
What does zustimmen mean? (\textquotedblleft{}To agree (with)\textquotedblright{}.) Say it once more.
\end{tcolorbox}
\section[die Zustimmung]{die Zustimmung — the approval, the consent}

\label{lesson:GE-C334-zustimmung}



\begin{quote}
Before the new one: say the German for to deny, then the German for to agree.
\end{quote}

\subsection*{You'll want to know: die Zustimmung}
\textbf{die Zustimmung} — \textquotedblleft{}the approval, the consent\textquotedblright{}.

\begin{quote}
Für den Ausflug brauchen die Kinder die Zustimmung der Eltern.
\end{quote}

\begin{quote}
\emph{For the trip the children need their parents' consent.}
\end{quote}

\emph{Zustimmung finden} — to meet with approval.

\begin{cousinweb}
\emph{Zustimmung} is the noun of \emph{zustimmen}, formed with \emph{-ung}; the verb is made from \emph{Stimme} \textquotedblleft{}voice, vote\textquotedblright{} (EtymDB).
\end{cousinweb}

\begin{grammarlens}[title={What you've built}]
One more word for telling, explaining and describing at B1.
\end{grammarlens}

\subsection*{Your turn}
\begin{itemize}
\raggedright
  \item \emph{Say it:} \emph{die Zustimmung}
  \item \emph{Say it:} \emph{die Zustimmung}, once more
  \item \emph{Say it:} a sentence of your own with \emph{die Zustimmung}
  \item \emph{Recall:} say the German for to deny, then the German for to agree, then say \emph{die Zustimmung} again
\end{itemize}

\begin{tcolorbox}[breakable,colback=teal!4,colframe=teal!35!black,title={Before you move on}]
What does die Zustimmung mean? (\textquotedblleft{}The approval, the consent\textquotedblright{}.) Say it once more.
\end{tcolorbox}
\section[einverstanden]{einverstanden — agreed, in agreement}

\label{lesson:GE-C334-einverstanden}



\begin{quote}
Before the new one: say the German for to agree, then the German for the approval.
\end{quote}

\subsection*{You'll want to know: einverstanden}
\textbf{einverstanden} — \textquotedblleft{}agreed, in agreement\textquotedblright{}.

\begin{quote}
Meine Eltern sind einverstanden: Ich darf am Samstag zur Party gehen.
\end{quote}

\begin{quote}
\emph{My parents agree: I am allowed to go to the party on Saturday.}
\end{quote}

\emph{Mit etwas einverstanden sein} — to agree to something. On its own, \emph{Einverstanden!} means OK, agreed.

\begin{cousinweb}
\emph{einverstanden} is the old participle of \emph{sich einverstehen} \textquotedblleft{}to reach an understanding\textquotedblright{}, built on \emph{verstehen} \textquotedblleft{}to understand\textquotedblright{}, from Old High German \emph{firstān} and Proto-Germanic \emph{frastāną} (EtymDB), and so on \emph{stehen} \textquotedblleft{}to stand\textquotedblright{}. EtymDB lists the parts, not the compound.
\end{cousinweb}

\begin{grammarlens}[title={What you've built}]
One more word for telling, explaining and describing at B1.
\end{grammarlens}

\subsection*{Your turn}
\begin{itemize}
\raggedright
  \item \emph{Say it:} \emph{einverstanden}
  \item \emph{Say it:} \emph{einverstanden}, once more
  \item \emph{Say it:} a sentence of your own with \emph{einverstanden}
  \item \emph{Recall:} say the German for to agree, then the German for the approval, then say \emph{einverstanden} again
\end{itemize}

\begin{tcolorbox}[breakable,colback=teal!4,colframe=teal!35!black,title={Before you move on}]
What does einverstanden mean? (\textquotedblleft{}Agreed, in agreement\textquotedblright{}.) Say it once more.
\end{tcolorbox}
\section[zugeben]{zugeben — to admit (a fault)}

\label{lesson:GE-C334-zugeben}



\begin{quote}
Before the new one: say the German for the approval, then the German for agreed.
\end{quote}

\subsection*{You'll want to know: zugeben}
\textbf{zugeben} — \textquotedblleft{}to admit (a fault)\textquotedblright{}.

\begin{quote}
Er gibt zu, dass er den Termin vergessen hat.
\end{quote}

\begin{quote}
\emph{He admits that he forgot the appointment.}
\end{quote}

Separable and strong: \emph{er gibt zu}, \emph{er hat zugegeben}. \emph{Zugegeben, …} — admittedly, …

\begin{cousinweb}
\emph{zugeben} is built from \emph{zu} and \emph{geben} \textquotedblleft{}to give\textquotedblright{}, inherited from Old High German \emph{geban} and Proto-Germanic \emph{gebaną} (EtymDB): you give a point to the other side.
\end{cousinweb}

\begin{grammarlens}[title={What you've built}]
One more word for telling, explaining and describing at B1.
\end{grammarlens}

\subsection*{Your turn}
\begin{itemize}
\raggedright
  \item \emph{Say it:} \emph{zugeben}
  \item \emph{Say it:} \emph{zugeben}, once more
  \item \emph{Say it:} a sentence of your own with \emph{zugeben}
  \item \emph{Recall:} say the German for the approval, then the German for agreed, then say \emph{zugeben} again
\end{itemize}

\begin{tcolorbox}[breakable,colback=teal!4,colframe=teal!35!black,title={Before you move on}]
What does zugeben mean? (\textquotedblleft{}To admit (a fault)\textquotedblright{}.) Say it once more.
\end{tcolorbox}
\section[einräumen]{einräumen — to concede; to put away (into a cupboard)}

\label{lesson:GE-C334-einraeumen}



\begin{quote}
Before the new one: say the German for agreed, then the German for to admit.
\end{quote}

\subsection*{You'll want to know: einräumen}
\textbf{einräumen} — \textquotedblleft{}to concede; to put away (into a cupboard)\textquotedblright{}.

\begin{quote}
Der Trainer räumt ein, dass seine Mannschaft schlecht gespielt hat.
\end{quote}

\begin{quote}
\emph{The coach concedes that his team played badly.}
\end{quote}

Separable: \emph{er räumt ein}, \emph{er hat eingeräumt}. Literally to give room to something: \emph{die Teller in den Schrank einräumen}.

\begin{cousinweb}
\emph{einräumen} is built on \emph{räumen} \textquotedblleft{}to make room\textquotedblright{}, inherited from Old High German \emph{rūmen} and Proto-Germanic \emph{rūmijaną} (EtymDB), from \emph{Raum} \textquotedblleft{}room\textquotedblright{}: you give the other side some room.
\end{cousinweb}

\begin{grammarlens}[title={What you've built}]
One more word for telling, explaining and describing at B1.
\end{grammarlens}

\subsection*{Your turn}
\begin{itemize}
\raggedright
  \item \emph{Say it:} \emph{einräumen}
  \item \emph{Say it:} \emph{einräumen}, once more
  \item \emph{Say it:} a sentence of your own with \emph{einräumen}
  \item \emph{Recall:} say the German for agreed, then the German for to admit, then say \emph{einräumen} again
\end{itemize}

\begin{tcolorbox}[breakable,colback=teal!4,colframe=teal!35!black,title={Before you move on}]
What does einräumen mean? (\textquotedblleft{}To concede; to put away (into a cupboard)\textquotedblright{}.) Say it once more.
\end{tcolorbox}
